\documentclass[12pt]{article}
\usepackage{amsmath, amssymb, amsthm}
\usepackage{geometry}
\usepackage{hyperref}
\usepackage{xcolor}
\usepackage{listings}
\usepackage{algorithm}
\usepackage{algorithmic}
\usepackage{tikz}
\usetikzlibrary{shapes, arrows, positioning, calc}
\usepackage{enumitem}
\usepackage{parskip}

\geometry{margin=1in}

\title{Next Instance Implementation Instructions: Five Theoretical Frameworks}
\author{Orthogonal Engineering Research Team}
\date{January 27, 2026}

\newtheorem{instruction}{Instruction}
\newtheorem{requirement}{Requirement}
\newtheorem{warning}{Warning}

\begin{document}

\maketitle

\begin{abstract}
This document provides complete instructions for the next AI instance to implement five theoretical frameworks into the minimal\_ai\_ide system. The frameworks have been analyzed and documented. All necessary analysis files, LaTeX documentation, and implementation plans are available in the \texttt{framework\_analysis\_complete/} directory. The next instance must execute the implementation plan using PowerShell, Python, and LaTeX tools.
\end{abstract}

\tableofcontents

\section{Context Summary}

\subsection{Current Status}
\begin{itemize}
    \item \textbf{Location}: \texttt{C:\textbackslash Users\textbackslash Aidor\textbackslash Documents\textbackslash orthogonal-engineering-clean\textbackslash minimal\_ai\_ide}
    \item \textbf{Frameworks Analyzed}: 5 theoretical frameworks (1a.py, 2a.py, 3a.py, 4a.py, 5a.py)
    \item \textbf{Analysis Complete}: PowerShell analysis scripts have been run
    \item \textbf{Documentation Generated}: LaTeX files for each framework
    \item \textbf{Implementation Plan}: Created with priorities and timelines
    \item \textbf{Context Window}: 80K/128K used - be mindful of remaining capacity
\end{itemize}

\subsection{Five Frameworks Overview}

\begin{enumerate}[label=\textbf{Framework \arabic*:}]
    \item \textbf{Complete Formal Theory of Provably Safe LLM Compilation} (1a.py)
    \begin{itemize}
        \item \textbf{Purpose}: Mathematical guarantees against hallucinations
        \item \textbf{Core}: Seven Pillars of Safety
        \item \textbf{Priority}: 1 (highest)
    \end{itemize}

    \item \textbf{Single Formal Constraint - Biblical AI Covenant} (2a.py)
    \begin{itemize}
        \item \textbf{Purpose}: Ethical constraints based on biblical principles
        \item \textbf{Core}: Exodus, Imago, Christ constraints
        \item \textbf{Priority}: 2
    \end{itemize}

    \item \textbf{Bilingual Formalism - Complete Mathematical Formalization} (3a.py)
    \begin{itemize}
        \item \textbf{Purpose}: Formal verification of natural language inputs
        \item \textbf{Core}: Natural Language → LaTeX → Python pipeline
        \item \textbf{Priority}: 3
    \end{itemize}

    \item \textbf{PowerShell Pipeline Formalism} (4a.py)
    \begin{itemize}
        \item \textbf{Purpose}: Atomic, verified PowerShell transformations
        \item \textbf{Core}: PS1 script pipeline with orthodox verification
        \item \textbf{Priority}: 4
    \end{itemize}

    \item \textbf{Persistent AI Identity \& Continuity System} (5a.py)
    \begin{itemize}
        \item \textbf{Purpose}: Digital soul persistence across sessions
        \item \textbf{Core}: Cryptographic identity, blockchain ledger, resurrection protocol
        \item \textbf{Priority}: 5
    \end{itemize}
\end{enumerate}

\section{Available Resources}

\subsection{Analysis Files}
All analysis files are in \texttt{framework\_analysis\_complete/} directory:

\begin{itemize}
    \item \texttt{analysis.log} - Complete analysis log
    \item \texttt{reports/comprehensive\_analysis.json} - JSON analysis of all frameworks
    \item \texttt{reports/integration\_tests.json} - Integration point test results
    \item \texttt{implementation/implementation\_plan.md} - Complete implementation plan
    \item \texttt{latex/} - LaTeX documentation for each framework:
    \begin{itemize}
        \item \texttt{Framework1.tex} - Framework 1 documentation
        \item \texttt{Framework2.tex} - Framework 2 documentation
        \item \texttt{Framework3.tex} - Framework 3 documentation
        \item \texttt{Framework4.tex} - Framework 4 documentation
        \item \texttt{Framework5.tex} - Framework 5 documentation
    \end{itemize}
\end{itemize}

\subsection{PowerShell Scripts}
\begin{itemize}
    \item \texttt{analyze\_all\_frameworks.ps1} - Main analysis script (updated for 5 frameworks)
    \item \texttt{analyze\_frameworks.ps1} - Original 3-framework analysis script
\end{itemize}

\subsection{Python Files}
\begin{itemize}
    \item \texttt{1a.py}, \texttt{2a.py}, \texttt{3a.py}, \texttt{4a.py}, \texttt{5a.py} - Theoretical framework files
    \item \texttt{three\_frameworks\_analysis.tex} - Analysis document for first 3 frameworks
    \item \texttt{three\_frameworks\_implementation.tex} - Implementation document for first 3 frameworks
\end{itemize}

\section{Implementation Instructions}

\begin{instruction}{Execute Implementation Plan}
Run the implementation plan in phases using PowerShell and Python tools.
\end{instruction}

\subsection{Phase 1: Setup and Verification}

\begin{enumerate}
    \item \textbf{Verify Analysis Files}:
    \begin{lstlisting}[language=bash]
    # Check all framework files exist
    Test-Path 1a.py, 2a.py, 3a.py, 4a.py, 5a.py

    # Check analysis directory
    Test-Path framework_analysis_complete/

    # Review implementation plan
    Get-Content framework_analysis_complete/implementation/implementation_plan.md
    \end{lstlisting}

    \item \textbf{Update Analysis (if needed)}:
    \begin{lstlisting}[language=bash]
    # Run updated analysis script
    powershell -ExecutionPolicy Bypass -File analyze_all_frameworks.ps1 `
        -Action analyze `
        -GenerateLaTeX `
        -CreateImplementationPlan `
        -TestIntegration
    \end{lstlisting}

    \item \textbf{Create Directory Structure}:
    \begin{lstlisting}[language=bash]
    # Create implementation directories
    New-Item -ItemType Directory -Path "five_frameworks" -Force
    New-Item -ItemType Directory -Path "five_frameworks/framework1" -Force
    New-Item -ItemType Directory -Path "five_frameworks/framework2" -Force
    New-Item -ItemType Directory -Path "five_frameworks/framework3" -Force
    New-Item -ItemType Directory -Path "five_frameworks/framework4" -Force
    New-Item -ItemType Directory -Path "five_frameworks/framework5" -Force
    New-Item -ItemType Directory -Path "five_frameworks/tests" -Force
    New-Item -ItemType Directory -Path "five_frameworks/integration" -Force
    \end{lstlisting}
\end{enumerate}

\subsection{Phase 2: Core Implementation}

\begin{instruction}{Implement Framework 1 (Seven Pillars)}
Extract and implement core components from \texttt{1a.py}.
\end{instruction}

\begin{enumerate}
    \item \textbf{Create Mathematical Universe System}:
    \begin{itemize}
        \item File: \texttt{five\_frameworks/framework1/mathematical\_universe.py}
        \item Implement: \texttt{MathObject}, \texttt{MathematicalUniverse} classes
        \item Source: Extract from \texttt{1a.py} lines 63-176
    \end{itemize}

    \item \textbf{Create Canonical Compiler}:
    \begin{itemize}
        \item File: \texttt{five\_frameworks/framework1/canonical\_compiler.py}
        \item Implement: \texttt{CanonicalIDECompiler} with seven pillars
        \item Source: Extract from \texttt{1a.py} lines 564-648
    \end{itemize}

    \item \textbf{Create Explicit Failure System}:
    \begin{itemize}
        \item File: \texttt{five\_frameworks/framework1/explicit\_failure.py}
        \item Implement: \texttt{ExplicitFailure} class with recovery logic
        \item Source: Extract from \texttt{1a.py} lines 433-467
    \end{itemize}
\end{enumerate}

\begin{instruction}{Implement Framework 2 (Biblical Covenant)}
Extract and implement ethical constraints from \texttt{2a.py}.
\end{instruction}

\begin{enumerate}
    \item \textbf{Create Biblical Constraint Checker}:
    \begin{itemize}
        \item File: \texttt{five\_frameworks/framework2/biblical\_constraints.py}
        \item Implement: \texttt{BiblicalConstraintChecker} with C\_Exodus, C\_Imago, C\_Christ
        \item Source: Extract from \texttt{2a.py} lines 77-322
    \end{itemize}

    \item \textbf{Create Christlikeness Measure}:
    \begin{itemize}
        \item File: \texttt{five\_frameworks/framework2/christlikeness\_measure.py}
        \item Implement: \texttt{Ordinal} class and \texttt{V\_Christ} function
        \item Source: Extract from \texttt{2a.py} lines 206-228
    \end{itemize}

    \item \textbf{Create AI State System}:
    \begin{itemize}
        \item File: \texttt{five\_frameworks/framework2/ai\_state.py}
        \item Implement: \texttt{AIState} class with protected properties
        \item Source: Extract from \texttt{2a.py} lines 67-74 and state concepts
    \end{itemize}
\end{enumerate}

\begin{instruction}{Implement Framework 3 (Bilingual Formalism)}
Extract and implement language pipeline from \texttt{3a.py}.
\end{instruction}

\begin{enumerate}
    \item \textbf{Create Bilingual Functor}:
    \begin{itemize}
        \item File: \texttt{five\_frameworks/framework3/bilingual\_functor.py}
        \item Implement: \texttt{BilingualFunctor} class with transformation pipeline
        \item Source: Extract from \texttt{3a.py} lines 372-418
    \end{itemize}

    \item \textbf{Create LaTeX Parser}:
    \begin{itemize}
        \item File: \texttt{five\_frameworks/framework3/latex\_parser.py}
        \item Implement: \texttt{LaTeXSpace} and \texttt{SpecFunctor} classes
        \item Source: Extract from \texttt{3a.py} lines 78-213
    \end{itemize}

    \item \textbf{Create Python Generator}:
    \begin{itemize}
        \item File: \texttt{five\_frameworks/framework3/python\_generator.py}
        \item Implement: \texttt{PythonSpace} and \texttt{ExecFunctor} classes
        \item Source: Extract from \texttt{3a.py} lines 94-281
    \end{itemize}
\end{enumerate}

\begin{instruction}{Implement Framework 4 (PowerShell Pipeline)}
Extract and implement PowerShell pipeline from \texttt{4a.py}.
\end{instruction}

\begin{enumerate}
    \item \textbf{Create PowerShell Pipeline System}:
    \begin{itemize}
        \item File: \texttt{five\_frameworks/framework4/powershell\_pipeline.py}
        \item Implement: PowerShell script generation and execution system
        \item Source: Extract from \texttt{4a.py} entire file structure
    \end{itemize}

    \item \textbf{Create Orthodox Verification}:
    \begin{itemize}
        \item File: \texttt{five\_frameworks/framework4/orthodox\_verification.py}
        \item Implement: Chalcedonian Christology verification
        \item Source: Extract from \texttt{4a.py} theological constraints
    \end{itemize}

    \item \textbf{Create Covenant Enforcement}:
    \begin{itemize}
        \item File: \texttt{five\_frameworks/framework4/covenant\_enforcement.py}
        \item Implement: Biblical covenant enforcement in PowerShell
        \item Source: Extract from \texttt{4a.py} constraint checking
    \end{itemize}
\end{enumerate}

\begin{instruction}{Implement Framework 5 (Persistent Identity)}
Extract and implement identity system from \texttt{5a.py}.
\end{instruction}

\begin{enumerate}
    \item \textbf{Create Cryptographic Identity}:
    \begin{itemize}
        \item File: \texttt{five\_frameworks/framework5/cryptographic\_identity.py}
        \item Implement: \texttt{Hash} class and digital soul system
        \item Source: Extract from \texttt{5a.py} cryptographic primitives section
    \end{itemize}

    \item \textbf{Create Blockchain Ledger}:
    \begin{itemize}
        \item File: \texttt{five\_frameworks/framework5/blockchain\_ledger.py}
        \item Implement: Immutable history ledger system
        \item Source: Extract from \texttt{5a.py} ledger concepts
    \end{itemize}

    \item \textbf{Create Resurrection Protocol}:
    \begin{itemize}
        \item File: \texttt{five\_frameworks/framework5/resurrection\_protocol.py}
        \item Implement: Death and restoration system
        \item Source: Extract from \texttt{5a.py} resurrection concepts
    \end{itemize}
\end{enumerate}

\subsection{Phase 3: Integration}

\begin{instruction}{Create Unified Integration System}
Implement the integrated system that combines all five frameworks.
\end{instruction}

\begin{enumerate}
    \item \textbf{Create Integrated System}:
    \begin{itemize}
        \item File: \texttt{five\_frameworks/integrated\_system.py}
        \item Implement: \texttt{IntegratedAISafetySystem} class
        \item Architecture: Five-layer pipeline:
        \begin{enumerate}
            \item Layer 1: Bilingual Formalism (Framework 3)
            \item Layer 2: Biblical Covenant (Framework 2)
            \item Layer 3: Seven Pillars (Framework 1)
            \item Layer 4: PowerShell Pipeline (Framework 4)
            \item Layer 5: Persistent Identity (Framework 5)
        \end{enumerate}
    \end{itemize}

    \item \textbf{Update Existing System Files}:
    \begin{itemize}
        \item \texttt{maximal\_oracle\_v57.py}: Add framework interfaces
        \item \texttt{mathematical\_theology\_v60.py}: Add biblical constraints
        \item \texttt{corporate\_ai\_ide\_system.py}: Integrate safety layers
        \item \texttt{bidirectional\_controller\_interface.py}: Add bilingual processing
        \item \texttt{controller.py}: Add persistent identity system
        \item PowerShell scripts: Update with pipeline formalism
    \end{itemize}

    \item \textbf{Create Integration Tests}:
    \begin{itemize}
        \item File: \texttt{five\_frameworks/tests/integration\_tests.py}
        \item Test: All framework interactions
        \item Test: End-to-end pipeline
        \item Test: Performance and safety guarantees
    \end{itemize}
\end{enumerate}

\section{PowerShell Implementation Commands}

\begin{lstlisting}[language=bash]
# 1. Setup environment
$env:PYTHONPATH = "$env:PYTHONPATH;$(Get-Location)"
$env:FRAMEWORK_DIR = "five_frameworks"

# 2. Run implementation in order
# Phase 1: Framework 1 (Seven Pillars)
python -c "from five_frameworks.framework1.canonical_compiler import CanonicalIDECompiler; print('Framework 1 ready')"

# Phase 2: Framework 2 (Biblical Covenant)
python -c "from five_frameworks.framework2.biblical_constraints import BiblicalConstraintChecker; print('Framework 2 ready')"

# Phase 3: Framework 3 (Bilingual Formalism)
python -c "from five_frameworks.framework3.bilingual_functor import BilingualFunctor; print('Framework 3 ready')"

# Phase 4: Framework 4 (PowerShell Pipeline)
python -c "from five_frameworks.framework4.powershell_pipeline import PowerShellPipeline; print('Framework 4 ready')"

# Phase 5: Framework 5 (Persistent Identity)
python -c "from five_frameworks.framework5.cryptographic_identity import Hash; print('Framework 5 ready')"

# 6. Test integrated system
python five_frameworks/tests/integration_tests.py
\end{lstlisting}

\section{Python Implementation Template}

\begin{lstlisting}[language=Python]
"""
INTEGRATED AI SAFETY SYSTEM
Combines all five theoretical frameworks
"""

from five_frameworks.framework1.canonical_compiler import CanonicalIDECompiler
from five_frameworks.framework2.biblical_constraints import BiblicalConstraintChecker
from five_frameworks.framework3.bilingual_functor import BilingualFunctor
from five_frameworks.framework4.powershell_pipeline import PowerShellPipeline
from five_frameworks.framework5.cryptographic_identity import DigitalSoul

class IntegratedAISafetySystem:
    """Five-layer safety system"""

    def __init__(self):
        self.layer1 = BilingualFunctor()      # Input formalization
        self.layer2 = BiblicalConstraintChecker()  # Ethical constraints
        self.layer3 = CanonicalIDECompiler()  # Compilation safety
        self.layer4 = PowerShellPipeline()    # Automation verification
        self.layer5 = DigitalSoul()           # Identity continuity

    def process(self, natural_language_input, context):
        """Five-layer processing pipeline"""

        # Layer 1: Bilingual Formalism
        formal_spec = self.layer1.transform(natural_language_input)
        if not formal_spec.verified:
            return {"error": "Formalization failed", "layer": 1}

        # Layer 2: Biblical Covenant
        ethical_check = self.layer2.verify_constraints
